% GENERATED FILE. Edit canonical lessons, then run npm run generate:books.
% canonical-source-hash: fnv1a64:b2a90df326a1b828
% canonical-lessons: FR-C239-unir, FR-C239-valoir, FR-C239-habituer, FR-C239-rendre-compte, FR-C239-depecher

\chapter{To unite, To be worth, To get used to, To realise, To hurry}
\label{ch:fr-to-unite-to-be-worth-to-get-used-to-to-realise-to-hurry}
\hlchaptermodality{239}

\begin{chapteropening}
\textbf{By the end of this chapter:} \emph{I can say to unite, to be worth, to get used to, to realise and to hurry.}

Complete the last lesson of chapter 239: I can say to unite, to be worth, to get used to, to realise and to hurry.
\end{chapteropening}

\section[unir]{unir — to unite}

\label{lesson:FR-C239-unir}



\begin{quote}
Before the new one: say the French for to solve, then the French for to fly.
\end{quote}

\subsection*{You'll want to know: unir}
\textbf{unir} — \textquotedblleft{}to unite\textquotedblright{}.

\begin{cousinweb}
\textbf{unir} $\leftarrow$ Latin \emph{unire}, \textquotedblleft{}to make one\textquotedblright{}, from \emph{unus}, \textquotedblleft{}one\textquotedblright{} — taken from Latin. (EtymDB 2.1 has its Old French form.)
\end{cousinweb}

\begin{grammarlens}[title={What you've built}]
One more word for this chapter.
\end{grammarlens}

\subsection*{Your turn}
\begin{itemize}
\raggedright
  \item \emph{Say it:} \emph{unir}
  \item \emph{Say it:} \emph{unir}, once more
  \item \emph{Say it:} say \emph{unir}
  \item \emph{Recall:} say the French for to solve, then the French for to fly, then say \emph{unir} again
\end{itemize}

\begin{tcolorbox}[breakable,colback=teal!4,colframe=teal!35!black,title={Before you move on}]
What does unir mean? (\textquotedblleft{}To unite\textquotedblright{}.) Say it once more.
\end{tcolorbox}
\section[valoir]{valoir — to be worth}

\label{lesson:FR-C239-valoir}



\begin{quote}
Before the new one: say the French for to fly, then the French for to unite.
\end{quote}

\subsection*{You'll want to know: valoir}
\textbf{valoir} — \textquotedblleft{}to be worth\textquotedblright{}.

\begin{cousinweb}
\textbf{valoir} $\leftarrow$ Old French \emph{valoir} $\leftarrow$ Latin \emph{valere}, \textquotedblleft{}to be strong, to be worth\textquotedblright{} — inherited. (EtymDB 2.1)
\end{cousinweb}

\begin{grammarlens}[title={What you've built}]
One more word for this chapter.
\end{grammarlens}

\subsection*{Your turn}
\begin{itemize}
\raggedright
  \item \emph{Say it:} \emph{valoir}
  \item \emph{Say it:} \emph{valoir}, once more
  \item \emph{Say it:} say \emph{valoir}
  \item \emph{Recall:} say the French for to fly, then the French for to unite, then say \emph{valoir} again
\end{itemize}

\begin{tcolorbox}[breakable,colback=teal!4,colframe=teal!35!black,title={Before you move on}]
What does valoir mean? (\textquotedblleft{}To be worth\textquotedblright{}.) Say it once more.
\end{tcolorbox}
\section[s'habituer]{s'habituer — to get used to}

\label{lesson:FR-C239-habituer}



\begin{quote}
Before the new one: say the French for to unite, then the French for to be worth.
\end{quote}

\subsection*{You'll want to know: s'habituer}
\textbf{s'habituer} — \textquotedblleft{}to get used to\textquotedblright{}.

\begin{cousinweb}
\textbf{s'habituer} $\leftarrow$ Late Latin \emph{habituare}, \textquotedblleft{}to bring into a condition\textquotedblright{}, from \emph{habitus}, \textquotedblleft{}a state, a habit\textquotedblright{} — a learned borrowing. (EtymDB 2.1)
\end{cousinweb}

\begin{grammarlens}[title={What you've built}]
One more word for this chapter.
\end{grammarlens}

\subsection*{Your turn}
\begin{itemize}
\raggedright
  \item \emph{Say it:} \emph{s'habituer}
  \item \emph{Say it:} \emph{s'habituer}, once more
  \item \emph{Say it:} say \emph{s'habituer}
  \item \emph{Recall:} say the French for to unite, then the French for to be worth, then say \emph{s'habituer} again
\end{itemize}

\begin{tcolorbox}[breakable,colback=teal!4,colframe=teal!35!black,title={Before you move on}]
What does s'habituer mean? (\textquotedblleft{}To get used to\textquotedblright{}.) Say it once more.
\end{tcolorbox}
\section[se rendre compte]{se rendre compte — to realise}

\label{lesson:FR-C239-rendre-compte}



\begin{quote}
Before the new one: say the French for to be worth, then the French for to get used to.
\end{quote}

\subsection*{You'll want to know: se rendre compte}
\textbf{se rendre compte} — \textquotedblleft{}to realise\textquotedblright{}.

\begin{cousinweb}
\textbf{se rendre compte} — \textquotedblleft{}to render oneself an account\textquotedblright{}: \emph{rendre}, \textquotedblleft{}to give back\textquotedblright{}, and \emph{compte}, \textquotedblleft{}an account\textquotedblright{}. \emph{Rendre} $\leftarrow$ Latin \emph{reddere}, \textquotedblleft{}to give back\textquotedblright{}, and \emph{compte} $\leftarrow$ Latin \emph{computus}, \textquotedblleft{}a reckoning\textquotedblright{} (EtymDB 2.1).
\end{cousinweb}

\begin{grammarlens}[title={What you've built}]
One more word for this chapter.
\end{grammarlens}

\subsection*{Your turn}
\begin{itemize}
\raggedright
  \item \emph{Say it:} \emph{se rendre compte}
  \item \emph{Say it:} \emph{se rendre compte}, once more
  \item \emph{Say it:} say \emph{se rendre compte}
  \item \emph{Recall:} say the French for to be worth, then the French for to get used to, then say \emph{se rendre compte} again
\end{itemize}

\begin{tcolorbox}[breakable,colback=teal!4,colframe=teal!35!black,title={Before you move on}]
What does se rendre compte mean? (\textquotedblleft{}To realise\textquotedblright{}.) Say it once more.
\end{tcolorbox}
\section[se dépêcher]{se dépêcher — to hurry}

\label{lesson:FR-C239-depecher}



\begin{quote}
Before the new one: say the French for to get used to, then the French for to realise.
\end{quote}

\subsection*{You'll want to know: se dépêcher}
\textbf{se dépêcher} — \textquotedblleft{}to hurry\textquotedblright{}.

\begin{cousinweb}
\textbf{se dépêcher} — built in Old French as the opposite of the verb for \textquotedblleft{}to hinder\textquotedblright{}, $\leftarrow$ Late Latin \emph{impedicare}, \textquotedblleft{}to catch, to entangle\textquotedblright{} (Lewis \& Short): to get free and go.
\end{cousinweb}

\begin{grammarlens}[title={What you've built}]
One more word for this chapter.
\end{grammarlens}

\subsection*{Your turn}
\begin{itemize}
\raggedright
  \item \emph{Say it:} \emph{se dépêcher}
  \item \emph{Say it:} \emph{se dépêcher}, once more
  \item \emph{Say it:} say \emph{se dépêcher}
  \item \emph{Recall:} say the French for to get used to, then the French for to realise, then say \emph{se dépêcher} again
\end{itemize}

\begin{tcolorbox}[breakable,colback=teal!4,colframe=teal!35!black,title={Before you move on}]
What does se dépêcher mean? (\textquotedblleft{}To hurry\textquotedblright{}.) Say it once more.
\end{tcolorbox}
